% !TeX root = ../main.tex
\section{Cạnh tranh với thị trường}

Hiện nay, người dùng có một số lựa chọn khi muốn mua hoặc trải nghiệm gaming gear, nhưng mỗi lựa chọn đều có hạn chế riêng.

\subsection{Mua trực tiếp tại shop}

Người dùng có thể đến cửa hàng để cầm thử sản phẩm. Tuy nhiên, trải nghiệm tại shop thường rất ngắn, không phản ánh đúng cảm giác khi chơi game trong nhiều giờ. Ngoài ra, không phải shop nào cũng có đủ mẫu để thử.

\subsection{Xem review trên YouTube, TikTok hoặc cộng đồng}

Review giúp người dùng tham khảo thông tin, nhưng trải nghiệm gear mang tính cá nhân rất cao. Một con chuột hợp với reviewer chưa chắc hợp với người dùng. Một loại switch được nhiều người khen chưa chắc phù hợp với thói quen gõ hoặc chơi game của từng cá nhân.

\subsection{Mua hàng online rồi đổi trả}

Một số sàn thương mại điện tử hoặc shop có chính sách đổi trả, nhưng quá trình đổi trả thường phức tạp. Người dùng vẫn phải thanh toán toàn bộ giá trị sản phẩm trước, đồng thời có thể gặp rủi ro nếu sản phẩm đã qua sử dụng hoặc không đáp ứng điều kiện đổi trả.

\subsection{Mua gear cũ}

Mua gear cũ giúp tiết kiệm chi phí, nhưng người dùng khó kiểm tra tình trạng thật của thiết bị. Ngoài ra, nếu mua nhầm, họ vẫn phải mất thời gian bán lại.

\subsection{Lợi thế cạnh tranh của nền tảng}

Nền tảng thuê/thử gear gaming có lợi thế vì tập trung trực tiếp vào nhu cầu \textbf{trải nghiệm thực tế trước khi mua}. Thay vì yêu cầu người dùng mua ngay, nền tảng cho phép họ trả một khoản phí nhỏ để dùng thử trong vài ngày.

So với các lựa chọn hiện có, sản phẩm có lợi thế ở các điểm:

\begin{itemize}
    \item Chi phí thử thấp hơn nhiều so với mua mới.
    \item Người dùng được trải nghiệm trong môi trường chơi game thật.
    \item Có thể thử nhiều mẫu trước khi quyết định mua.
    \item Có đánh giá thực tế từ cộng đồng người đã thuê.
    \item Shop có thêm kênh tiếp cận khách hàng.
    \item Có thể kết hợp mô hình thuê, bán và try-before-buy trong cùng một nền tảng.
\end{itemize}
